\documentclass[11pt]{article}

\usepackage[margin=1in]{geometry}
\usepackage{amsmath,amssymb}
\usepackage{booktabs}
\usepackage{graphicx}
\usepackage{float}
\usepackage{xcolor}
\usepackage{microtype}
\usepackage[numbers,sort&compress]{natbib}
\usepackage[font=small,labelfont=bf]{caption}
\PassOptionsToPackage{hyphens}{url}
\usepackage[colorlinks=true,linkcolor=blue!50!black,citecolor=blue!50!black,urlcolor=blue!50!black]{hyperref}
\Urlmuskip=0mu plus 1mu\relax
\emergencystretch=1.2em
\renewcommand{\bibfont}{\scriptsize\raggedright}
\setlength{\bibsep}{1pt}
\usepackage[compact]{titlesec}
\setlength{\textfloatsep}{14pt plus 2pt minus 4pt}

\newif\ifdarkmode
\ifdefined\DARKMODE \darkmodetrue \else \darkmodefalse \fi
\ifdarkmode
  \input{darkstyle}
  \graphicspath{{figures/dark/}}
\else
  \graphicspath{{figures/}}
\fi

\newcommand{\singleauthorinfo}[1]{\\[0.25em]{\small #1}}
\newcommand{\q}[1]{\texttt{#1}}
% Generated by paper3/scripts/build_claims.py. DO NOT EDIT.
% Every value comes from paper3/derived/claims.json.
% Unsupported cache-hit/decode-I/O quantities intentionally have no macros.
% Hardware and GGUF artifact inventory
\newcommand{\PthreeGpuVramMiB}{12282}
\newcommand{\PthreeGpuMarketedVramGB}{12}
\newcommand{\PthreeGpuPowerCapWatts}{80}
\newcommand{\PthreeCpuRamGiB}{62}
\newcommand{\PthreeKernelVersion}{7.0.0-29-generic}
\newcommand{\PthreeNvidiaDriverVersion}{580.167.08}
\newcommand{\PthreeCudaVersion}{13.0}
\newcommand{\PthreeSwapTotalGiB}{8}
\newcommand{\PthreeSwapUsedBaselineGiB}{7.3}
\newcommand{\PthreeFlashFileCount}{3}
\newcommand{\PthreeFlashFileBytes}{72546461344}
\newcommand{\PthreeFlashFileGiB}{67.564}
\newcommand{\PthreeFlashLogicalElements}{176943899520}
\newcommand{\PthreeFlashEffectiveBpw}{3.28}
\newcommand{\PthreeFlashLayers}{48}
\newcommand{\PthreeFlashExperts}{512}
\newcommand{\PthreeFlashActiveExperts}{10}
\newcommand{\PthreeFlashRoutedTensorCount}{144}
\newcommand{\PthreeFlashRoutedBytes}{39845888000}
\newcommand{\PthreeFlashRoutedGiB}{37.109}
\newcommand{\PthreeFlashRoutedLogicalElements}{120795955200}
\newcommand{\PthreeFlashMeanRoutedLayerGiB}{0.773}
\newcommand{\PthreeFlashMeanExpertLayerSliceMiB}{1.546}
\newcommand{\PthreeFlashNoReuseActiveExpertMiBPerToken}{742.188}
\newcommand{\PthreeFlashNoReuseActiveExpertGiBPerTwoFiveSixTokens}{185.547}
\newcommand{\PthreeGptOssFileCount}{1}
\newcommand{\PthreeGptOssFileBytes}{63387346208}
\newcommand{\PthreeGptOssFileGiB}{59.034}
\newcommand{\PthreeGptOssLogicalElements}{116829156672}
\newcommand{\PthreeGptOssEffectiveBpw}{4.341}
\newcommand{\PthreeGptOssLayers}{36}
\newcommand{\PthreeGptOssExperts}{128}
\newcommand{\PthreeGptOssActiveExperts}{4}
\newcommand{\PthreeGptOssRoutedTensorCount}{216}
\newcommand{\PthreeGptOssRoutedBytes}{61073326080}
\newcommand{\PthreeGptOssRoutedGiB}{56.879}
\newcommand{\PthreeGptOssRoutedLogicalElements}{114701598720}
\newcommand{\PthreeGptOssMeanRoutedLayerGiB}{1.58}
\newcommand{\PthreeGptOssMeanExpertLayerSliceMiB}{12.64}
\newcommand{\PthreeGptOssNoReuseActiveExpertMiBPerToken}{1820.127}
\newcommand{\PthreeGptOssNoReuseActiveExpertGiBPerTwoFiveSixTokens}{455.032}
\newcommand{\PthreeFlashNonRoutedTableBytes}{28800138240}
\newcommand{\PthreeFlashNonRoutedTableGiB}{26.822}
\newcommand{\PthreeExpertSliceSizeRatio}{8.17}
% Three-run generation sequences and full-invocation device counters
\newcommand{\PthreeFlashInitialN}{3}
\newcommand{\PthreeFlashInitialGenMean}{9.93}
\newcommand{\PthreeFlashInitialReadMeanGiB}{103.3}
\newcommand{\PthreeFlashInitialPromptOne}{4.9}
\newcommand{\PthreeFlashInitialGenOne}{9.4}
\newcommand{\PthreeFlashInitialReadOneGiB}{104.2}
\newcommand{\PthreeFlashInitialCachedPostOneGiB}{20}
\newcommand{\PthreeFlashInitialMemAvailablePreOneGiB}{21.6}
\newcommand{\PthreeFlashInitialPromptTwo}{5}
\newcommand{\PthreeFlashInitialGenTwo}{10.4}
\newcommand{\PthreeFlashInitialReadTwoGiB}{102.1}
\newcommand{\PthreeFlashInitialCachedPostTwoGiB}{19.9}
\newcommand{\PthreeFlashInitialMemAvailablePreTwoGiB}{21.7}
\newcommand{\PthreeFlashInitialPromptThree}{4.9}
\newcommand{\PthreeFlashInitialGenThree}{10}
\newcommand{\PthreeFlashInitialReadThreeGiB}{103.6}
\newcommand{\PthreeFlashInitialCachedPostThreeGiB}{19.9}
\newcommand{\PthreeFlashInitialMemAvailablePreThreeGiB}{21.8}
\newcommand{\PthreeFlashQuietN}{3}
\newcommand{\PthreeFlashQuietGenMean}{14.33}
\newcommand{\PthreeFlashQuietReadMeanGiB}{79.47}
\newcommand{\PthreeFlashQuietPromptOne}{9.2}
\newcommand{\PthreeFlashQuietGenOne}{14.7}
\newcommand{\PthreeFlashQuietReadOneGiB}{84.3}
\newcommand{\PthreeFlashQuietMemAvailablePostOneGiB}{56}
\newcommand{\PthreeFlashQuietPromptTwo}{12}
\newcommand{\PthreeFlashQuietGenTwo}{14.4}
\newcommand{\PthreeFlashQuietReadTwoGiB}{82.1}
\newcommand{\PthreeFlashQuietMemAvailablePostTwoGiB}{55}
\newcommand{\PthreeFlashQuietPromptThree}{3.3}
\newcommand{\PthreeFlashQuietGenThree}{13.9}
\newcommand{\PthreeFlashQuietReadThreeGiB}{72}
\newcommand{\PthreeFlashQuietMemAvailablePostThreeGiB}{57}
\newcommand{\PthreeGptOssInitialN}{3}
\newcommand{\PthreeGptOssInitialGenMean}{9.07}
\newcommand{\PthreeGptOssInitialReadMeanGiB}{123.7}
\newcommand{\PthreeGptOssInitialPromptOne}{3.5}
\newcommand{\PthreeGptOssInitialGenOne}{9.3}
\newcommand{\PthreeGptOssInitialReadOneGiB}{126.2}
\newcommand{\PthreeGptOssInitialMemAvailablePostOneGiB}{31}
\newcommand{\PthreeGptOssInitialPromptTwo}{3.3}
\newcommand{\PthreeGptOssInitialGenTwo}{8.7}
\newcommand{\PthreeGptOssInitialReadTwoGiB}{124.1}
\newcommand{\PthreeGptOssInitialMemAvailablePostTwoGiB}{30}
\newcommand{\PthreeGptOssInitialPromptThree}{3.4}
\newcommand{\PthreeGptOssInitialGenThree}{9.2}
\newcommand{\PthreeGptOssInitialReadThreeGiB}{120.8}
\newcommand{\PthreeGptOssInitialMemAvailablePostThreeGiB}{30}
\newcommand{\PthreeFlashQuietMinusInitialMean}{4.4}
\newcommand{\PthreeFlashQuietRelativeToInitialPct}{44.3}
% Non-randomized Flash expert-placement observations
\newcommand{\PthreeFlashMoeFortyFourCpuLayers}{44}
\newcommand{\PthreeFlashMoeFortyFourNonCpuExpertLayers}{4}
\newcommand{\PthreeFlashMoeFortyFourPrompt}{5.2}
\newcommand{\PthreeFlashMoeFortyFourGen}{10.6}
\newcommand{\PthreeFlashMoeFortyFourReadGiB}{101.4}
\newcommand{\PthreeFlashMoeFortyFourDescriptiveChangePct}{6.7}
\newcommand{\PthreeFlashMoeFortyTwoCpuLayers}{42}
\newcommand{\PthreeFlashMoeFortyTwoNonCpuExpertLayers}{6}
\newcommand{\PthreeFlashMoeFortyTwoPrompt}{5.5}
\newcommand{\PthreeFlashMoeFortyTwoGen}{10.6}
\newcommand{\PthreeFlashMoeFortyTwoReadGiB}{99.8}
\newcommand{\PthreeFlashMoeFortyTwoDescriptiveChangePct}{6.7}
% Single ascending prompt-scale sweeps
\newcommand{\PthreeFlashPrefillN}{3}
\newcommand{\PthreeFlashPrefillContext}{8192}
\newcommand{\PthreeFlashPrefillApproxTokensOne}{160}
\newcommand{\PthreeFlashPrefillPromptOne}{6.6}
\newcommand{\PthreeFlashPrefillGenOne}{3.7}
\newcommand{\PthreeFlashPrefillRequestedDecodeOne}{16}
\newcommand{\PthreeFlashPrefillApproxTokensTwo}{640}
\newcommand{\PthreeFlashPrefillPromptTwo}{13.1}
\newcommand{\PthreeFlashPrefillGenTwo}{3.7}
\newcommand{\PthreeFlashPrefillRequestedDecodeTwo}{16}
\newcommand{\PthreeFlashPrefillApproxTokensThree}{2560}
\newcommand{\PthreeFlashPrefillPromptThree}{20.1}
\newcommand{\PthreeFlashPrefillGenThree}{3.2}
\newcommand{\PthreeFlashPrefillRequestedDecodeThree}{16}
\newcommand{\PthreeGptOssPrefillN}{3}
\newcommand{\PthreeGptOssPrefillContext}{8192}
\newcommand{\PthreeGptOssPrefillApproxTokensOne}{160}
\newcommand{\PthreeGptOssPrefillPromptOne}{4.2}
\newcommand{\PthreeGptOssPrefillGenOne}{5.4}
\newcommand{\PthreeGptOssPrefillRequestedDecodeOne}{64}
\newcommand{\PthreeGptOssPrefillReadOneGiB}{117.9}
\newcommand{\PthreeGptOssPrefillApproxTokensTwo}{640}
\newcommand{\PthreeGptOssPrefillPromptTwo}{9.1}
\newcommand{\PthreeGptOssPrefillGenTwo}{4.9}
\newcommand{\PthreeGptOssPrefillRequestedDecodeTwo}{64}
\newcommand{\PthreeGptOssPrefillReadTwoGiB}{127}
\newcommand{\PthreeGptOssPrefillApproxTokensThree}{2560}
\newcommand{\PthreeGptOssPrefillPromptThree}{23.5}
\newcommand{\PthreeGptOssPrefillGenThree}{5.8}
\newcommand{\PthreeGptOssPrefillRequestedDecodeThree}{64}
\newcommand{\PthreeGptOssPrefillReadThreeGiB}{138.5}
% Row-backed quality cells, plus explicitly separate aggregate cells
\newcommand{\PthreeFlashMathRawN}{25}
\newcommand{\PthreeFlashMathRawCorrect}{18}
\newcommand{\PthreeFlashMathRawPct}{72}
\newcommand{\PthreeFlashMathRawWilsonLowPct}{52.4}
\newcommand{\PthreeFlashMathRawWilsonHighPct}{85.7}
\newcommand{\PthreeFlashMathRawErrors}{1}
\newcommand{\PthreeFlashMathRawTruncations}{0}
\newcommand{\PthreeFlashCodeSavedN}{17}
\newcommand{\PthreeFlashCodeSavedCorrect}{16}
\newcommand{\PthreeFlashCodeSavedPct}{94.1}
\newcommand{\PthreeFlashCodeSavedWilsonLowPct}{73}
\newcommand{\PthreeFlashCodeSavedWilsonHighPct}{99}
\newcommand{\PthreeFlashCodeSavedErrors}{0}
\newcommand{\PthreeFlashCodeSavedTruncations}{0}
\newcommand{\PthreeFlashGsmN}{50}
\newcommand{\PthreeFlashGsmCorrect}{48}
\newcommand{\PthreeFlashGsmPct}{96}
\newcommand{\PthreeFlashGsmWilsonLowPct}{86.5}
\newcommand{\PthreeFlashGsmWilsonHighPct}{98.9}
\newcommand{\PthreeFlashGsmErrors}{0}
\newcommand{\PthreeFlashGsmTruncations}{0}
\newcommand{\PthreeFlashMathLenientN}{25}
\newcommand{\PthreeFlashMathLenientCorrect}{23}
\newcommand{\PthreeFlashMathLenientPct}{92}
\newcommand{\PthreeFlashMathLenientWilsonLowPct}{75}
\newcommand{\PthreeFlashMathLenientWilsonHighPct}{97.8}
\newcommand{\PthreeFlashMathFormatOnlyUpgrades}{5}
\newcommand{\PthreeFlashCodeAggregateN}{20}
\newcommand{\PthreeFlashCodeAggregateCorrect}{17}
\newcommand{\PthreeFlashCodeAggregatePct}{85}
\newcommand{\PthreeFlashCodeAggregateWilsonLowPct}{64}
\newcommand{\PthreeFlashCodeAggregateWilsonHighPct}{94.8}
\newcommand{\PthreeGptOssMathRawN}{25}
\newcommand{\PthreeGptOssMathRawCorrect}{12}
\newcommand{\PthreeGptOssMathRawPct}{48}
\newcommand{\PthreeGptOssMathRawWilsonLowPct}{30}
\newcommand{\PthreeGptOssMathRawWilsonHighPct}{66.5}
\newcommand{\PthreeGptOssMathRawErrors}{1}
\newcommand{\PthreeGptOssMathRawTruncations}{0}
\newcommand{\PthreeGptOssCodeSavedN}{20}
\newcommand{\PthreeGptOssCodeSavedCorrect}{18}
\newcommand{\PthreeGptOssCodeSavedPct}{90}
\newcommand{\PthreeGptOssCodeSavedWilsonLowPct}{69.9}
\newcommand{\PthreeGptOssCodeSavedWilsonHighPct}{97.2}
\newcommand{\PthreeGptOssCodeSavedErrors}{0}
\newcommand{\PthreeGptOssCodeSavedTruncations}{0}
\newcommand{\PthreeGptOssGsmN}{50}
\newcommand{\PthreeGptOssGsmCorrect}{46}
\newcommand{\PthreeGptOssGsmPct}{92}
\newcommand{\PthreeGptOssGsmWilsonLowPct}{81.2}
\newcommand{\PthreeGptOssGsmWilsonHighPct}{96.8}
\newcommand{\PthreeGptOssGsmErrors}{0}
\newcommand{\PthreeGptOssGsmTruncations}{0}
\newcommand{\PthreeGptOssMathLenientN}{25}
\newcommand{\PthreeGptOssMathLenientCorrect}{20}
\newcommand{\PthreeGptOssMathLenientPct}{80}
\newcommand{\PthreeGptOssMathLenientWilsonLowPct}{60.9}
\newcommand{\PthreeGptOssMathLenientWilsonHighPct}{91.1}
\newcommand{\PthreeGptOssMathFormatOnlyUpgrades}{8}
\newcommand{\PthreeGptOssCodeAggregateN}{20}
\newcommand{\PthreeGptOssCodeAggregateCorrect}{18}
\newcommand{\PthreeGptOssCodeAggregatePct}{90}
\newcommand{\PthreeGptOssCodeAggregateWilsonLowPct}{69.9}
\newcommand{\PthreeGptOssCodeAggregateWilsonHighPct}{97.2}
\newcommand{\PthreeFlashCodeMissingRows}{3}
% Shared harness protocol
\newcommand{\PthreeGenerationContext}{4096}
\newcommand{\PthreeGenerationRequestedDecode}{256}
\newcommand{\PthreeQuietGateMinMemAvailableGiB}{35}
\newcommand{\PthreeQuietGateMaxLoadExclusive}{3}
\newcommand{\PthreeQuietGateConsecutiveChecks}{10}
\newcommand{\PthreeQuietGateCheckIntervalSeconds}{30}
\newcommand{\PthreeQuietGateDurationMinutes}{5}
\newcommand{\PthreeQualityContext}{8192}
\newcommand{\PthreeQualityReasoningEffort}{medium}
\newcommand{\PthreeQualityTemperature}{1}
\newcommand{\PthreeQualityTopP}{0.95}
\newcommand{\PthreeQualityTopK}{20}
\newcommand{\PthreeQualityConcurrency}{1}
\newcommand{\PthreeMathMaxTokens}{8192}
\newcommand{\PthreeCodeMaxTokens}{4096}
\newcommand{\PthreeGsmMaxTokens}{4096}
\newcommand{\PthreeGsmLimit}{50}
\newcommand{\PthreeDatasetSelectionSeed}{42}
\newcommand{\PthreeMathMatchedItems}{25}
\newcommand{\PthreeGsmMatchedItems}{50}
\newcommand{\PthreeHumanEvalFrozenItems}{20}


\title{\textbf{llama.cpp Throughput for Qwen3.8-Flash-Next\\
and gpt-oss-120b on a 62 GiB RAM Laptop}}

\hypersetup{
  pdftitle={llama.cpp Throughput for Qwen3.8-Flash-Next and gpt-oss-120b on a 62 GiB RAM Laptop},
  pdfauthor={Matthew Schwartz},
  pdfsubject={Decode-throughput measurements for two large mixture-of-experts GGUF deployments using llama.cpp's ordinary mmap path}
}

\author{Matthew Schwartz\singleauthorinfo{Independent researcher \enspace
\href{https://orcid.org/0009-0009-4171-7247}{ORCID 0009-0009-4171-7247}
\enspace \texttt{matthew.schwartz.research@gmail.com}}}
\date{September 2026}

\begin{document}
\maketitle

\begin{abstract}
Large mixture-of-experts (MoE) models can be deployed across system memory,
GPU memory, and storage, but file size alone does not predict generation
speed. We measured two public GGUF deployments with unmodified llama.cpp on
an Intel Core i9-14900HX laptop with \PthreeCpuRamGiB\,GiB of RAM and an
\PthreeGpuPowerCapWatts\,W RTX 4080 Laptop GPU. The engine used its ordinary
memory-mapped file path, with routed experts assigned to the CPU and GPU
offload requested for other supported layers. The \PthreeFlashFileGiB\,GiB
Qwen3.8-Flash-Next artifact decoded at
\PthreeFlashInitialGenOne--\PthreeFlashQuietGenOne\ tokens/s across two
sequential three-run groups under different recorded host states. A
\PthreeGptOssFileGiB\,GiB gpt-oss-120b artifact decoded at
\PthreeGptOssInitialGenTwo--\PthreeGptOssInitialGenOne\ tokens/s across three
runs. Each invocation requested \PthreeGenerationRequestedDecode\ output
tokens from one fixed prompt. These are decode rates: they exclude model
loading and prompt processing and do not establish interactive response
latency. The report provides run-level measurements, exact deployment
identities, and separate small task checks as a reference for reproducing
these configurations. Sequential run groups and whole-device storage
counters leave the causes of the observed rates unresolved. The contribution
is a hardware-specific throughput benchmark; no inference method was changed.
\end{abstract}

\section{Introduction}

A local deployment must deliver answers within a useful time budget as well
as fit the available memory. With 59--68\,GiB model files on a
\PthreeCpuRamGiB\,GiB RAM laptop, the operating system, inference buffers,
and weights compete for system memory, while a
\PthreeGpuMarketedVramGB\,GB GPU holds only part of the deployment.
Measuring generation speed is one step toward deciding whether such a
configuration can serve the intended workload.

Linux memory mapping lets an engine address file-backed pages without first
copying the whole artifact into RAM \citep{linuxmmap2026,linuxmm2026}.
llama.cpp users have already reported large, storage-backed MoE deployments
through this path \citep{llamacppdiscussion18758,llamacppdiscussion23324}.
The open practical question for a new configuration is its measured speed,
rather than whether memory mapping permits execution at all.

Mixture-of-experts models make this performance question especially relevant.
They store many feed-forward experts but route each token through only a few
\citep{shazeer2017moe,fedus2021switch,dai2024deepseekmoe}, so stored size and
active computation are not the same quantity. Purpose-built runtimes exploit
this structure through expert placement, prefetching, compression, and caching
\citep{edgemoe2023,mixtraloffload2023,hobbit2024,slicemoe2025,dymoe2026,
pagedweight2026}. The ordinary mmap path in llama.cpp provides a reference
configuration for evaluating these techniques \citep{gerganov2023llamacpp}.

We ask how fast Qwen3.8-Flash-Next and gpt-oss-120b decoded on this laptop
with a fixed llama.cpp checkout and CPU/GPU placement. ``Unmodified'' refers
to the source checkout; the command lines use explicit placement flags.

The \PthreeFlashFileGiB\,GiB Flash deployment completed generation at
\PthreeFlashInitialGenOne--\PthreeFlashQuietGenOne\ decode tokens/s across its
recorded host states, and the \PthreeGptOssFileGiB\,GiB gpt-oss deployment ran
at \PthreeGptOssInitialGenTwo--\PthreeGptOssInitialGenOne\ tokens/s. The
contribution is this benchmark and the configuration needed to repeat it.
Tensor inventories document the files, and small separate task samples
provide limited checks of the deployed models' answers. Prompt processing is
reported separately because decode speed alone does not determine response
latency.

The larger file exceeded installed system RAM, but not combined nominal RAM
and VRAM; the study did not measure simultaneous residency of every tensor
byte. The measurements therefore do not establish execution beyond combined
physical memory or identify a paging mechanism.
Section~\ref{sec:methods} defines the exact deployments and measurements,
Section~\ref{sec:results} reports their observed rates and supporting audits,
and the remaining sections explain how the baseline can and cannot be used.

\section{Methods}
\label{sec:methods}

\subsection{Machine, engine, and placement}

The host was a Linux laptop with an Intel Core i9-14900HX,
\PthreeCpuRamGiB\,GiB of
dual-channel DDR5 memory, a 2\,TB WD PC SN740 NVMe drive, and an NVIDIA RTX
4080 Laptop GPU with \PthreeGpuVramMiB\,MiB of VRAM at an
\PthreeGpuPowerCapWatts\,W power limit. The recorded baseline inventory lists
Linux \q{\PthreeKernelVersion}, NVIDIA driver
\q{\PthreeNvidiaDriverVersion}, and driver-reported CUDA compatibility
\q{\PthreeCudaVersion}. It also
records \PthreeSwapTotalGiB\,GiB of swap, with
\PthreeSwapUsedBaselineGiB\,GiB in use at that snapshot. These are inventory
values, not per-run measurements of software state, memory pressure, or swap
traffic. Run logs
record llama.cpp build \q{b1-035e227}; the associated clean source checkout
currently resolves to
\q{035e2273}\allowbreak\q{1a7fd70b}\allowbreak\q{9854b3a2}\allowbreak
\q{d64ec68e}\allowbreak\q{9b1a45d3}. It was compiled with CUDA for
compute capability 8.9. The surviving build record identifies GNU C++ 13.3.0
and CUDA toolkit compiler 12.6.85; the release artifact records the remaining
toolchain and CMake options. For the Flash runs, that revision was a commit in the
then-unmerged Qwen4Exp support pull request; the pull request subsequently
merged on August 27, 2026 \citep{llamacpp2026qwen4exppr}. The gpt-oss path did
not depend on that experimental architecture support.

Both command lines used llama.cpp's default mmap behavior, requested GPU
offload with \q{-ngl 99}, and assigned routed MoE tensors to the CPU path with
\q{--n-cpu-moe 999}. Thus the requested placement combines CPU expert
execution with GPU execution of other supported layers. The exact flags and
engine revision define the baseline.

\subsection{Artifacts and tensor inventory}

GGUF is the model container used by llama.cpp. The first subject was the
three-shard Unsloth
\path{Qwen3.8-Flash-Next-UD-IQ1_S} artifact at Hugging Face revision
\q{d3bc75ee}\allowbreak\q{6ccef3ef}\allowbreak\q{c1e228ec}\allowbreak
\q{00a6cc2c}\allowbreak\q{db1e2249}. Its GGUF metadata declares
\PthreeFlashLayers\ blocks, \PthreeFlashExperts\ experts, and
\PthreeFlashActiveExperts\ active experts per token
\citep{qwenflash2026,unsloth2026flashnext}. The second was
\path{gpt-oss-120b-MXFP4.gguf} at revision
\q{238abdd2}\allowbreak\q{90bb874b}\allowbreak\q{90a5da1b}\allowbreak
\q{4549881b}\allowbreak\q{7d05c091}, with \PthreeGptOssLayers\
blocks, \PthreeGptOssExperts\ experts, and \PthreeGptOssActiveExperts\ active
per token \citep{openai2025gptoss,ggmlorg2025gptossgguf}.

We parsed the frozen GGUF headers and summed stored bytes and logical tensor
elements by standardized tensor name. Effective whole-file bits per enumerated
logical tensor element (bpw) is
$8\times$ artifact bytes divided by enumerated logical tensor elements. The
inventory is descriptive of the deployed files; it does not measure which
pages the engine read during a token.

\begin{table}[H]
\centering
\footnotesize
\begin{tabular}{lrrrrl}
\toprule
Artifact & File GiB & File bpw & Routed experts GiB & Experts/active & Notable other tensor \\
\midrule
Flash UD-IQ1\_S & \PthreeFlashFileGiB & \PthreeFlashEffectiveBpw &
\PthreeFlashRoutedGiB & \PthreeFlashExperts/\PthreeFlashActiveExperts &
\PthreeFlashNonRoutedTableGiB\ GiB lookup table \\
gpt-oss MXFP4 & \PthreeGptOssFileGiB & \PthreeGptOssEffectiveBpw &
\PthreeGptOssRoutedGiB & \PthreeGptOssExperts/\PthreeGptOssActiveExperts & --- \\
\bottomrule
\end{tabular}
\caption{Deployed-file inventory from GGUF headers. Flash's
\q{per\_layer\_token\_embd.weight} lookup table is outside the
routed-expert tensor set. File bpw includes all tensors and container overhead.}
\label{tab:inventory}
\end{table}

\subsection{Timing protocol and storage counters}

Each generation invocation used one fixed technical-explanation prompt, a
\PthreeGenerationContext-token context limit, a request for
\PthreeGenerationRequestedDecode\ output tokens, and llama.cpp's
\q{-st} timing output. The harness saved the reported prompt and generation
throughput but not the generated text for these timing runs. It did not set a
fixed seed or force greedy decoding. The exact prompt was: \emph{Write a
detailed technical explanation of how an operating system page cache interacts
with memory-mapped files during sequential and random access patterns.}

The Flash campaign contains an initial three-run sequence labeled as a loaded
host and a later three-run sequence fired by a quiet gate after available
memory remained at least \PthreeQuietGateMinMemAvailableGiB\,GiB and one-minute
load remained below \PthreeQuietGateMaxLoadExclusive\ for
\PthreeQuietGateDurationMinutes\ minutes. The gpt-oss campaign contains one
initial three-run sequence. These
are sequential observations, not randomized or interleaved condition samples.
The host-state labels describe the campaign; they do not isolate cache size or
background load as a cause.

Immediately before process launch and after process exit, the harness sampled
the read-sector counter for the physical NVMe device in
\q{/proc/diskstats}. We call the difference a \emph{whole-invocation
device-read delta}. It includes every phase of \q{llama-cli} and reads by
other processes on the same device. The generation rate, by contrast, is
llama.cpp's decode timing. These two intervals must not be divided to obtain
decode bytes per token, cache misses, or effective decode bandwidth. Memory
fields also differ across groups: the initial Flash rows record post-run
\q{Cached} and pre-run \q{MemAvailable}, while later groups record post-run
\q{MemAvailable}. We therefore do not put them on one cache-size axis.

\subsection{Quality checks and alternate scoring}

Separate saved evaluations used frozen subsets constructed with dataset seed
\PthreeDatasetSelectionSeed. The surviving files contain the same
\PthreeMathMatchedItems\ MATH and \PthreeGsmMatchedItems\ GSM8K item IDs for
both models. Both campaigns targeted the same
\PthreeHumanEvalFrozenItems-item HumanEval+ subset, although the Flash file is
incomplete as described below. The GSM8K subset consists of the first
\PthreeGsmMatchedItems\ entries of a seed-42 sample of 100 test items sorted
by source index, not a direct random sample of 50. Servers used an \PthreeQualityContext-token
context, temperature \PthreeQualityTemperature, top-$p$
\PthreeQualityTopP, top-$k$ \PthreeQualityTopK, requested reasoning effort
\q{\PthreeQualityReasoningEffort}, and concurrency
\PthreeQualityConcurrency. Maximum requested output was
\PthreeMathMaxTokens\ tokens for MATH, \PthreeCodeMaxTokens\ for HumanEval+,
and \PthreeGsmMaxTokens\ for GSM8K
\citep{hendrycks2021math,chen2021humaneval,liu2023evalplus,
cobbe2021gsm8k}. The harness sent one user message per item and no separate
system message. MATH appended ``End your response with: Answer: <final
answer>'' to the problem; GSM8K appended ``End your response with: Answer:
<number>'' to the question. HumanEval+ asked for the complete runnable Python
function, including its signature and imports, in one Python code block with
nothing afterward, then supplied the task's function prompt. MATH is reported
both under the original harness scorer and a
more lenient post-hoc rescore. The original scorer already extracts an answer,
normalizes common LaTeX and fraction forms, and applies numeric tolerance. The
post-hoc rescore accepts additional presentation variants and simple numeric
equivalents. The original score remains primary; the two scores are not
merged. One retained MATH row from each model records a timeout error and was
counted incorrect.
The harness could retry a failed request automatically; the saved rows do not
record attempt counts. Task scores describe retained outcomes, not necessarily
one generation attempt per item.

The saved Flash HumanEval+ file contains only \PthreeFlashCodeSavedN\ rows
(\PthreeFlashCodeSavedCorrect\ correct), while the aggregate ledger reports
\PthreeFlashCodeAggregateCorrect/\PthreeFlashCodeAggregateN. Because the
\PthreeFlashCodeMissingRows\ missing item rows cannot be
audited, we omit the aggregate Flash code score from the scientific result and
retain the discrepancy in the artifact. Wilson 95\% intervals accompany the
row-backed proportions \citep{wilson1927}.

\section{Results}
\label{sec:results}

\subsection{Observed decode throughput}
\vspace{0.15em}

The \PthreeFlashFileGiB\,GiB Flash deployment produced
\PthreeFlashInitialGenOne--\PthreeFlashQuietGenOne\ decode tokens/s across its
six recorded invocations. The \PthreeGptOssFileGiB\,GiB gpt-oss deployment
produced \PthreeGptOssInitialGenTwo--\PthreeGptOssInitialGenOne\ tokens/s
across three invocations. Every principal invocation reached a generation
timing result. Table~\ref{tab:generation} gives the individual rates and the
separately scoped storage counters. The ranges are observed extrema, not
confidence intervals or estimates of run-to-run variability.

Prompt processing was slower than decode in these runs. The initial Flash
group reported \PthreeFlashInitialPromptOne, \PthreeFlashInitialPromptTwo,
and \PthreeFlashInitialPromptThree\ input tokens/s; the later group reported
\PthreeFlashQuietPromptOne, \PthreeFlashQuietPromptTwo, and
\PthreeFlashQuietPromptThree. The gpt-oss group reported
\PthreeGptOssInitialPromptOne, \PthreeGptOssInitialPromptTwo, and
\PthreeGptOssInitialPromptThree\ input tokens/s. Loading time and prompt
processing must therefore be measured alongside decode before using these
configurations for latency-sensitive work. The single prompt does not support
a prediction for longer contexts.

\begin{table}[H]
\centering
\small
\begin{tabular}{llrrl}
\toprule
Artifact & Recorded group & $n$ & Decode tok/s & Whole-invocation device reads (GiB) \\
\midrule
Flash & initial loaded label & \PthreeFlashInitialN &
\PthreeFlashInitialGenOne\ / \PthreeFlashInitialGenTwo\ / \PthreeFlashInitialGenThree &
\PthreeFlashInitialReadOneGiB\ / \PthreeFlashInitialReadTwoGiB\ / \PthreeFlashInitialReadThreeGiB \\
Flash & quiet gate & \PthreeFlashQuietN &
\PthreeFlashQuietGenOne\ / \PthreeFlashQuietGenTwo\ / \PthreeFlashQuietGenThree &
\PthreeFlashQuietReadOneGiB\ / \PthreeFlashQuietReadTwoGiB\ / \PthreeFlashQuietReadThreeGiB \\
gpt-oss & initial loaded label & \PthreeGptOssInitialN &
\PthreeGptOssInitialGenOne\ / \PthreeGptOssInitialGenTwo\ / \PthreeGptOssInitialGenThree &
\PthreeGptOssInitialReadOneGiB\ / \PthreeGptOssInitialReadTwoGiB\ / \PthreeGptOssInitialReadThreeGiB \\
\bottomrule
\end{tabular}
\caption{Sequential generation observations. The read column is the physical
device's counter change over each complete process invocation, not expert
traffic or decode-only I/O. The observations are descriptive because the runs
were not randomized and host state changed between groups.}
\label{tab:generation}
\end{table}

The higher Flash rates in the quiet-gated sequence are an observed association
with a different host state, not an estimate of a cache-size treatment effect.
Likewise, the declining device-read deltas within that sequence do not prove
cache warming: the counter includes unrelated reads and its interval differs
from the decode timer. The data establish the observed throughput of the two
deployments, not why one sequence was faster.

\subsection{Stored tensors and the limits of I/O accounting}

Flash contains \PthreeFlashRoutedGiB\,GiB of tensors whose standardized names
identify routed experts. Its \PthreeFlashNonRoutedTableGiB\,GiB
\q{per\_layer\_token\_embd.weight} table is separate and does not match the
routed \q{*\_exps} naming rule. The audit does not infer a more specific
architectural role from that name alone. The mean stored
routed-expert slice is \PthreeFlashMeanExpertLayerSliceMiB\,MiB per
expert-layer for Flash and \PthreeGptOssMeanExpertLayerSliceMiB\,MiB for
gpt-oss, an \PthreeExpertSliceSizeRatio-fold storage difference.

From those inventories, reading every active routed-expert slice once with no
reuse would amount arithmetically to
\PthreeFlashNoReuseActiveExpertGiBPerTwoFiveSixTokens\,GiB per
\PthreeGenerationRequestedDecode\ Flash tokens and
\PthreeGptOssNoReuseActiveExpertGiBPerTwoFiveSixTokens\,GiB per
\PthreeGenerationRequestedDecode\ gpt-oss tokens. These are counterfactual storage sums, not
measured I/O. Because the recorded device deltas span additional tensors and
process phases, comparing them with the counterfactual does not yield a
page-cache hit rate.

\subsection{Sampled task performance}

\begin{table}[H]
\centering
\footnotesize
\begin{tabular}{lllrr}
\toprule
Artifact & Task and scoring rule & Correct/$n$ & Accuracy & Wilson 95\% interval \\
\midrule
Flash & MATH, original harness &
\PthreeFlashMathRawCorrect/\PthreeFlashMathRawN & \PthreeFlashMathRawPct\% &
\PthreeFlashMathRawWilsonLowPct--\PthreeFlashMathRawWilsonHighPct\% \\
Flash & MATH, post-hoc lenient &
\PthreeFlashMathLenientCorrect/\PthreeFlashMathLenientN & \PthreeFlashMathLenientPct\% &
\PthreeFlashMathLenientWilsonLowPct--\PthreeFlashMathLenientWilsonHighPct\% \\
Flash & GSM8K & \PthreeFlashGsmCorrect/\PthreeFlashGsmN & \PthreeFlashGsmPct\% &
\PthreeFlashGsmWilsonLowPct--\PthreeFlashGsmWilsonHighPct\% \\
gpt-oss & MATH, original harness &
\PthreeGptOssMathRawCorrect/\PthreeGptOssMathRawN & \PthreeGptOssMathRawPct\% &
\PthreeGptOssMathRawWilsonLowPct--\PthreeGptOssMathRawWilsonHighPct\% \\
gpt-oss & MATH, post-hoc lenient &
\PthreeGptOssMathLenientCorrect/\PthreeGptOssMathLenientN & \PthreeGptOssMathLenientPct\% &
\PthreeGptOssMathLenientWilsonLowPct--\PthreeGptOssMathLenientWilsonHighPct\% \\
gpt-oss & HumanEval+ &
\PthreeGptOssCodeSavedCorrect/\PthreeGptOssCodeSavedN & \PthreeGptOssCodeSavedPct\% &
\PthreeGptOssCodeSavedWilsonLowPct--\PthreeGptOssCodeSavedWilsonHighPct\% \\
gpt-oss & GSM8K & \PthreeGptOssGsmCorrect/\PthreeGptOssGsmN & \PthreeGptOssGsmPct\% &
\PthreeGptOssGsmWilsonLowPct--\PthreeGptOssGsmWilsonHighPct\% \\
\bottomrule
\end{tabular}
\caption{Row-backed task checks for the deployed configurations. The Flash
HumanEval+ aggregate is omitted because \PthreeFlashCodeMissingRows\ item rows
are missing. The two
MATH rows score the same saved responses under different disclosed rules.}
\label{tab:quality}
\end{table}

These samples show that the deployed configurations returned correct answers
on several tasks while operating through the measured stacks. The MATH result
depends materially on the disclosed scoring rule, so the original and post-hoc
values remain separate. These samples do not show that mmap preserves quality,
because neither artifact was also evaluated in a fully memory-resident
configuration and the artifacts are already quantized.

\section{Discussion}

\subsection{What the throughput baseline is useful for}

The measured decode rates give users of these exact files a starting point
for a local trial. They also specify a baseline that can be rerun under matched
conditions when evaluating an expert cache or offloader. A useful comparison
would keep the model file, prompt, generation settings, placement, and host
state fixed, and report loading, prompt processing, decode, and total response
time separately. The slow prompt rates here make that distinction practical:
a tolerable stream of output tokens can follow a long wait for the first one.

\subsection{What controls performance remains open}

The recorded groups show that throughput varied with host state, but their
measurement windows cannot assign that variation to cache residency,
background load, storage traffic, or another cause. Explaining decode traffic
would require phase-scoped and preferably process-attributed I/O, consistent
before/after residency metrics, fixed decoding, and randomized repeated host
states. Direct router or tensor-access instrumentation would then be needed to
distinguish expert popularity, lookup-table traffic, and reuse.

\section{Limitations and threats to validity}

This is one machine, one NVMe device, one engine commit, two artifacts, and
three sequential observations per principal group. The Flash artifact exceeds
physical RAM and gpt-oss is smaller; the paper does not cover models several
times larger than host memory. Host activity, thermal
state, page residency, and invocation order were not randomized or held fixed.
The timing harness retained extracted throughput lines rather than full CLI
outputs, and it did not set a fixed seed.

Device reads are whole-invocation, physical-device-global counters. They cannot
separate engine I/O from unrelated host reads. They also do not distinguish
loading, prefill, decode, and teardown. The memory fields differ across groups.
No comparison with a purpose-built offloader, a direct-I/O or memory-locking
configuration, or a fully memory-resident run of the same artifact was made.

Quality samples are small. MATH conclusions depend strongly on whether the
original harness score or more lenient post-hoc rescore is used, so both are
reported. Flash HumanEval+
has an unresolved aggregate/item-row mismatch and is omitted. Model and
hardware scope is exact: results should not be generalized to other MoEs,
quantizations, storage devices, operating systems, or engine revisions without
replication.

\section{Related work}

EdgeMoE, Fast MoE Offloading, HOBBIT, SliceMoE, DyMoE, and PagedWeight represent
an engineered lineage of expert placement, caching, prefetching, and mixed
precision \citep{edgemoe2023,mixtraloffload2023,hobbit2024,slicemoe2025,
dymoe2026,pagedweight2026}. A proposed llama.cpp \q{O\_DIRECT} expert streamer
targets pools far beyond RAM and reports improvements over mmap
\citep{llamacpp2026streampr}, and SSD-LLaMA builds a three-tier SSD, RAM, and
VRAM expert hierarchy into llama.cpp to run trillion-parameter MoEs on a
consumer PC \citep{liang2026ssdllama}. Those systems pursue speed, scale, or predictable
resource control. The measurements here document an unmodified mmap
configuration that could be reused in a future matched comparison; no such
comparison is made here.

An opt-in GPU-resident LRU cache for host-offloaded experts is proposed in
llama.cpp pull request 27861 \citep{llamacpp2026moecachepr}, still an open
draft when checked on September 28, 2026. Its author reports 18.4 versus 24.2
decode tokens/s in a matched cache-disabled/cache-enabled test with a larger
Flash quantization on a dual-Xeon, two-RTX-3090 system. A community port of the
same cache with an added admission gate reports 5.6 versus 10.0 decode tokens/s
for Flash IQ3\_XXS on a 12\,GB RTX 3060 with 64\,GB of RAM, in interleaved
arms on one binary \citep{llamacpp2026moecachepr}. These within-system
comparisons test an optimization. Its speedup cannot be transferred to the present
laptop, which used another artifact, placement, and engine revision. Testing
the cache here would require a new matched experiment.

A later community fork reports 55.54 decode tokens/s without speculation and
61.68 with multi-token prediction for Flash UD-Q3\_K\_XL on a
62.14\,GiB RAM desktop with a 16\,GB RTX 5070 Ti
\citep{generelschwerz2026flashreproduction}. Its September 14 runs use expert
caching, lazy loading of the per-layer embedding table, and other engine
changes. These are author-reported short-prompt results, not independently
reproduced measurements here. Similar RAM capacity does not make the CPU,
GPU, quantization, placement, prompt, or software comparable. The present
report preserves an older, pinned baseline; it does not recommend that
revision over newer runtimes or estimate their speedup on this laptop.

Community reports already establish storage-backed mmap execution
\citep{llamacppdiscussion18758,llamacppdiscussion23324}; the contribution here
is the auditable configuration and measurements for these exact deployments.
Recent deployment benchmarks also address adjacent questions: Alfarizy et al.\
compare a smaller MoE with dense models on consumer and edge devices
\citep{alfarizy2026consumer}, while a community Flash report sweeps placement
and loading modes on one RTX PRO 6000 workstation with 96\,GB system RAM
\citep{lukallm2026flashbenchmark}. Neither is a matched comparison with the
two laptop deployments reported here.

\section{Artifact availability and reproducibility}

The reproducibility artifact accompanying this report contains
checksummed, sanitized timing measurements and item-level quality outcomes,
the claims data, environment and artifact manifests, model revisions, and a
deterministic analysis-only checker. Without
executing a model, the checker verifies the exported files, recomputes reported
generation rates and row-backed task scores, and checks the reported artifact
byte lengths, effective bpw, and routed-expert totals. It does not include or
parse the private working logs or local GGUF headers. Full replication requires
downloading the pinned upstream artifacts and running the documented commands
on compatible hardware; the weights are linked rather than redistributed.
The version-1.0.0 artifact is identified by
\href{https://doi.org/10.5281/zenodo.22314465}{doi:10.5281/zenodo.22314465}.
Its code, sanitized evidence, and replication instructions share the immutable
\href{https://github.com/matthematics1137/research-artifacts/tree/moe-nvme-streaming-v1.0.0/papers/2026-moe-nvme-streaming}{versioned GitHub artifact path}.
The DOI identifies the frozen reproducibility bundle; it is not a separate
manuscript DOI or a claim of peer review.

\section{Conclusion}

On a \PthreeCpuRamGiB\,GiB RAM, \PthreeGpuMarketedVramGB\,GB VRAM laptop,
unmodified llama.cpp produced
\PthreeFlashInitialGenOne--\PthreeFlashQuietGenOne\ decode tokens/s with the
\PthreeFlashFileGiB\,GiB Qwen3.8-Flash-Next artifact and
\PthreeGptOssInitialGenTwo--\PthreeGptOssInitialGenOne\ tokens/s with the
\PthreeGptOssFileGiB\,GiB gpt-oss-120b artifact. The archived invocations and
configuration provide a baseline for repeating these deployments and testing
changes under matched conditions. Slow prompt processing also shows why
decode throughput must be separated from response latency. The observations
characterize these files on this machine; their short sequential design does
not explain the speed differences or establish quality preservation.

\medskip
\noindent\textbf{AI assistance.} Claude (Anthropic) assisted with evaluation
harness development and experiment orchestration. Claude and Codex (OpenAI)
assisted with evidence auditing, analysis and figure code, citation checking,
and manuscript drafting and revision. Matthew Schwartz directed the research
and is responsible for its methods, results, interpretation, claims,
citations, rights, and released artifacts.

\bibliographystyle{unsrtnat}
\bibliography{references}

\end{document}
